\chapter{Software Implementation}
\label{ch:implementation}

This chapter describes the modular C++ software structure developed for
reading robot states, accessing model information, and executing the
impedance controller in a reproducible manner, addressing Objective~3
of the thesis registration.

\section{Software Architecture}
\label{sec:sw_arch}

The software is structured in three layers, illustrated in
\cref{fig:sw_arch}:

\begin{enumerate}
  \item \textbf{Communication layer}: establishes and maintains the
        1\,kHz \abbr{FCI} connection via \texttt{libfranka}.
  \item \textbf{State and model layer}: reads joint positions, velocities,
        and the Jacobian; retrieves gravity and Coriolis vectors from the
        robot model.
  \item \textbf{Control layer}: evaluates the impedance law and returns
        joint torques within the 1\,ms deadline.
\end{enumerate}

\begin{figure}[htb]
\centering
\begin{tikzpicture}[
  box/.style={draw, rounded corners=3pt, minimum width=4.2cm,
               minimum height=0.9cm, align=center, font=\small},
  arr/.style={-{Latex}, thick}
]
  \node[box, fill=blue!10]  (hw)    at (0, 0.0) {Franka Panda\\(Hardware)};
  \node[box, fill=orange!15](fci)   at (0,-1.8) {\abbr{FCI} / \texttt{libfranka}};
  \node[box, fill=green!10] (state) at (0,-3.6) {State \& Model Layer\\($\mathbf{T}_e$, $\mathbf{J}$, $\bm{\tau}_g$)};
  \node[box, fill=red!10]   (ctrl)  at (0,-5.4) {Control Layer\\$\bm{\tau}=\mathbf{J}^\top\mathbf{F}$};
  \draw[arr] (hw)    -- (fci)   node[midway,right,font=\tiny]{1\,kHz};
  \draw[arr] (fci)   -- (state) node[midway,right,font=\tiny]{robot state};
  \draw[arr] (state) -- (ctrl)  node[midway,right,font=\tiny]{$\vq,\mathbf{T}_e,\mathbf{J}$};
  \draw[arr] (ctrl.east) -- ++(1.5,0) |- (fci.east)
    node[pos=0.25,right,font=\tiny]{$\bm{\tau}_{\mathrm{cmd}}$};
\end{tikzpicture}
\caption{Three-layer software architecture. The control layer evaluates
         the impedance law and returns joint torques to the FCI
         within the 1\,ms deadline.}
\label{fig:sw_arch}
\end{figure}

\section{State Acquisition and Model Access}
\label{sec:state_model}

Each 1\,ms control cycle extracts the following from
\texttt{franka::RobotState}:

\begin{itemize}
  \item \texttt{state.q}/\texttt{state.dq}: $\vq$ and $\vqd$.
  \item \texttt{state.O\_T\_EE}: $\mathbf{T}_e\in \SE$
        as a $4\times4$ column-major array of 16 doubles.
  \item \texttt{model.zeroJacobian(...)}: the $6\times7$ geometric
        Jacobian $\mathbf{J}(\vq)$ in the base frame.
  \item \texttt{state.gravity}: the $7\times1$ gravity vector
        $\bm{\tau}_g(\vq)$.
  \item \texttt{state.coriolis}: the $7\times1$ Coriolis vector
        $\bm{\tau}_c(\vq,\vqd)$.
\end{itemize}

\section{Control Loop Implementation}
\label{sec:control_impl}

The complete impedance control loop is implemented as a \texttt{libfranka}
torque callback.
The callback receives the current robot state and must return a
\texttt{franka::Torques} object within 1\,ms.
The following listing shows the full implementation as deployed in the
experiments:

\begin{lstlisting}[style=cppstyle,
  caption={Complete Cartesian impedance control callback},
  label={lst:full_ctrl}]
robot.control([&](const franka::RobotState& state,
    franka::Duration) -> franka::Torques {

  // -- State acquisition --
  Eigen::Map<const Eigen::Matrix<double,7,1>> q(state.q.data());
  Eigen::Map<const Eigen::Matrix<double,7,1>> dq(state.dq.data());

  // -- End-effector pose T_e from state.O_T_EE --
  Eigen::Affine3d T_e(Eigen::Matrix4d::Map(state.O_T_EE.data()));
  Eigen::Matrix3d R_e = T_e.rotation();
  Eigen::Vector3d p_e = T_e.translation();

  // -- Jacobian J(q) from model --
  std::array<double,42> jac_arr =
      model.zeroJacobian(franka::Frame::kEndEffector, state);
  Eigen::Map<const Eigen::Matrix<double,6,7>> J(jac_arr.data());

  // -- Cartesian velocity --
  Eigen::Matrix<double,6,1> xdot = J * dq;
  Eigen::Vector3d pdot  = xdot.head<3>();
  Eigen::Vector3d omega = xdot.tail<3>();

  // -- Pose error via T_err = T_d^{-1} T_e --
  Eigen::Matrix3d R_d = T_des.rotation();
  Eigen::Vector3d p_d = T_des.translation();
  Eigen::Vector3d e_p = R_d.transpose() * (p_e - p_d);

  Eigen::Matrix3d dR = R_d.transpose() * R_e;
  double phi = std::acos(std::clamp(0.5*(dR.trace()-1.0),-1.0,1.0));
  Eigen::Vector3d e_R = Eigen::Vector3d::Zero();
  if (phi > 1e-6)
      e_R = (phi/(2.0*std::sin(phi))) *
            Eigen::Vector3d(dR(2,1)-dR(1,2),
                            dR(0,2)-dR(2,0),
                            dR(1,0)-dR(0,1));

  // -- Impedance wrench (full SPD K and D matrices) --
  Eigen::Vector3d f = Kp * e_p - Dp * pdot;
  Eigen::Vector3d m = KR * e_R - DR * omega;
  Eigen::Matrix<double,6,1> F;
  F.head<3>() = f;  F.tail<3>() = m;

  // -- Joint torques: tau = J^T F --
  Eigen::Matrix<double,7,1> tau = J.transpose() * F;

  // -- Gravity and Coriolis compensation --
  Eigen::Map<const Eigen::Matrix<double,7,1>>
      tau_g(state.gravity.data()),
      tau_c(state.coriolis.data());
  tau += tau_g + tau_c;

  // -- Safety clipping --
  const double tau_max[] = {87,87,87,87,12,12,12};
  std::array<double,7> out{};
  for (size_t i=0; i<7; ++i)
    out[i] = std::max(-tau_max[i], std::min(tau_max[i], tau(i)));
  return franka::Torques(out);
});
\end{lstlisting}

\section{Safety and Reproducibility}
\label{sec:safety}

All experiments are conducted following these safety procedures:

\begin{itemize}
  \item Controller is first tested with $k_{\min} = 10\,\mathrm{N/m}$
        before increasing to operating stiffness.
  \item Joint torques are clipped to manufacturer limits at all times.
  \item The emergency stop is held by an observer throughout each trial.
  \item All stiffness matrices are verified to be \abbr{SPD} before controller
        activation (Cholesky factorisation check).
  \item Each experimental trial is repeated three times; results are
        averaged across repetitions to assess reproducibility.
\end{itemize}
